\newpage
\section{Eksploracyjna analiza danych}\label{sec:opis_danych}

Dane źródłowe pochodzą ze zbioru opublikowanego na konkursie podczas konferencji International Conference on Prognostics and Health Management (PHM08). Celem tego konkursu było wykorzystanie danych ze wskaźników silnikowych do oszacowania pozostałej liczby cykli dla danej floty silników.  Dane te nie zostały zebrane poprzez odczyt fizycznie zamontowanych urządzeń pomiarowych, lecz wygenerowane w~wyniku symulacji narzędzia C-MAPSS (Commercial Modular Aero-Propulsion System Simulation). C-MAPSS pozwala na symulację pracy silnika turbowentylatorowego klasy 90.000 lbf (funt-siła) ciągu (tej samej klasy, które zasilają Boeing-777). Symulacja deterioracji silnika została przeprowadzono poprzez sterowanie zmianą sprawności komponentów, przepływów oraz innych parametrów silnika. Dostęp do rzeczywistych danych - a~nie symulowanych - nie jest możliwy ze względu na wysokie koszty pozyskania połączone z~ochroną konkurencyjności i~praw własnościowych do danych \cite{saxena2008damage}.

\begin{longtable}{|p{1cm}||p{7cm}|p{2.2cm}|p{2cm}|}
\caption{\textbf{Zestawienie parametrów silnika}} \label{tab:parametry_silnika} \\
\toprule
 & Nazwa & Oznaczenie & Jednostka \\
L.p. &  &  &  \\
\midrule
\endfirsthead
\caption[]{Zestawienie parametrów silnika} \\
\toprule
 & Nazwa & Oznaczenie & Jednostka \\
L.p. &  &  &  \\
\midrule
\endhead
\midrule
\multicolumn{4}{r}{C.d. na następnej stronie} \\
\midrule
\endfoot
\bottomrule
\endlastfoot
0 & Temperatura całkowita na wlocie wentylatora & T2 & °R \\
1 & Temperatura całkowita na wylocie Sprężarki Niskiego Ciśnienia & T24 & °R \\
2 & Temperatura całkowita na wylocie Sprężarki Wysokiego Ciśnienia & T30 & °R \\
3 & Temperatura całkowita na wylocie Turbiny Niskiego Ciśnienia & T50 & °R \\
4 & Ciśnienie na wlocie wentylatora & P2 & psia \\
5 & Ciśnienie całkowite w~kanale zewnętrznym & P15 & psia \\
6 & Ciśnienie całkowite na wylocie Sprężarki Wysokiego Ciśnienia & P30 & psia \\
7 & Prędkość obrotowa wentylatora & Nf & obr/min \\
8 & Prędkość obrotowa wytwornicy spalin & Nc & obr/min \\
9 & Spręż silnika & epr & -- \\
10 & Ciśnienie statyczne na wylocie Sprężarki Wysokiego Ciśnienia & Ps30 & psia \\
11 & Stosunek wydatku paliwa do ciśnienia statycznego na wylocie Sprężarki Wysokiego Ciśnienia & phi & pps/psi \\
12 & Poprawiona prędkość obrotowa wentylatora & NRf & obr/min \\
13 & Poprawiona prędkość obrotowa wytwornicy spalin & NRc & obr/min \\
14 & Stosunek dwuprzepływowości & BPR & -- \\
15 & Stosunek wydatku paliwa do powietrza & farB & -- \\
16 & Entalpia wydatku masowego & htBleed & -- \\
17 & Wymagana prędkość obrotowa wentylatora & Nf\_dmd & obr/min \\
18 & Wymagana poprawiona prędkość obrotowa wentylatora & PCNfR\_dmd & obr/min \\
19 & Wydatek masowy chłodzenia Turbiny Wysokiego Ciśnienia & W31 & lbm/s \\
20 & Wydatek masowy chłodzenia Turbiny Niskiego Ciśnienia & W32 & lbm/s \\
\end{longtable}

gdzie:
\textbf{$^{\circ}R$} – stopnie Rankine’a (jednostka temperatury bezwzględnej w układzie anglosaskim), \textbf{psia} – funty na cal kwadratowy, wartość bezwzględna ciśnienia (ang. \textit{pounds per square inch absolute}), \textbf{obr~/~min} – obroty na minutę, \textbf{lbm~/~s} – funty masy na sekundę, miara ilości przepływającej masy (ang. \textit{pounds mass per second}), \textbf{pps~/~psi} – funty paliwa na sekundę na funt na cal kwadratowy (ang. \textit{pounds per second / pound per square inch}).

Przeprowadzono kompleksową analizę danych począwszy od analizy rozkładu maksymalnego przeżycia dla każdego z~silników. Utworzono histogram RUL (rysunek~\ref{fig:hist_max_life}), którego oś pionowa określa liczbę elementów z danego zakresu. Wykres wskazuje na niesymetryczny rozkład maksymalnej liczby przelatanych cykli. Histogram gromadzi się głównie wokół wartości 200 cykli, chociaż wyjątkowo trwałe silniki o~maksymalnym życiu większym niż 350 również znajdują się w~zbiorze danych.

\begin{figure}[H]
    \centering
    \caption{\textbf {Histogram Rozkładu Maksymalnej liczby cykli}}
    \includegraphics[width=0.8\linewidth]{Output/EDA/Histograms/Max Life_histogram.png}
    \label{fig:hist_max_life}
    \caption*{Źródło: opracowanie własne} 
\end{figure}

Ze zbioru danych usunięto następujące kolumny, które były w~całości puste lub miały tylko jedną wartość: \textit{Operational setting 3, T2, P2, epr, farB, $Nf_{dmd}$, $PCNfR_{dmd}$}
Wartości te są stałe ze względu na naturę symulacji: temperatura (\textit{T2}) i~ciśnienie na wlocie (\textit{P2}) są stałe, ponieważ symulowane są przeloty w~atmosferze ISA (standardowe wartości parametrów otoczenia). Stosunek ciśnień (\textit{epr}, spręż) silnika jest również stały, ponieważ zmiany ciśnienia na wylocie turbiny są relatywnie małe. To samo dotyczy wymaganej prędkości wentylatora (\textit{$Nf_{dmd}$}) oraz wydatku paliwa (\textit{farB}), ze względu na stałe obciążenie wartości te również pozostają bez zmian.

Obliczono RUL (ang. Remaining Useful Life), czyli pozostałą liczbę cykli do końca życia silnika i~wykonano analizę korelacji Pearsona, otrzymując rysunek~\ref{fig:pearson}.
Na podstawie macierzy korelacji stworzono dendrogram cech (rysunek \ref{fig:dendrogram}) w~celu ich zgrupowania. 
Przy progu odcięcia odległości korelacji równym 0.25 otrzymano grupy opisane w~tabeli \ref{tab:clusters}.
Za pomocą pakietu XGBoost wykonano analizę istotności cech (ang. \textit{Feature Importance}) uzyskując wykres słupkowy istotności cech (rysunek \ref{fig:xgb_importance}).

\begin{figure}[H]
    \centering
    \caption{\textbf{Mapa cieplna korelacji Pearsona cech zbioru danych}}
    \includegraphics[width=0.8\linewidth]{Output/Pearson Correlation Heatmap.png}
    \caption*{Źródło: opracowanie własne}
    \label{fig:pearson}
\end{figure}

\begin{figure}[H]
    \centering
    \caption{\textbf{Analiza klasteryzacji - dendrogram}\\Oś pionowa przedstawia odległość korelacji Pearsona między cechami.}
    \includegraphics[width=0.9\linewidth]{Output/Dendrogram.png} 
    \caption*{Źródło: opracowanie własne}
    \label{fig:dendrogram}
\end{figure}

\begin{table}[H]
\centering
\caption{\textbf{Wynik klasteryzacji}}
\label{tab:clusters}
\resizebox{\textwidth}{!}{
\centering
\begin{tabular}{lll}
\toprule
 &  & Opis \\
Grupa & Symbol &  \\
\midrule
\multirow[t]{5}{*}{1} & P30 & Ciśnienie całkowite na wylocie Sprężarki Wysokiego Ciśnienia \\
 & RUL & [-] \\
 & W31 & Wydatek masowy chłodzenia Turbiny Wysokiego Ciśnienia \\
 & W32 & Wydatek masowy chłodzenia Turbiny Niskiego Ciśnienia \\
 & phi & Stosunek wydatku paliwa do ciśnienia statycznego na wylocie Sprężarki Wysokiego Ciśnienia \\
\cline{1-3}
2 & Operational setting 1 & [-] \\
\cline{1-3}
3 & Operational setting 2 & [-] \\
\cline{1-3}
\multirow[t]{8}{*}{4} & BPR & Stosunek dwuprzepływowości \\
 & NRf & Poprawiona prędkość obrotowa wentylatora \\
 & Nf & Prędkość obrotowa wentylatora \\
 & Ps30 & Ciśnienie statyczne na wylocie Sprężarki Wysokiego Ciśnienia \\
 & T24 & Temperatura całkowita na wylocie Sprężarki Niskiego Ciśnienia \\
 & T30 & Temperatura całkowita na wylocie Sprężarki Wysokiego Ciśnienia \\
 & T50 & Temperatura całkowita na wylocie Turbiny Niskiego Ciśnienia \\
 & htBleed & Entalpia wydatku masowego \\
\cline{1-3}
5 & P15 & Ciśnienie całkowite w kanale zewnętrznym \\
\cline{1-3}
\multirow[t]{2}{*}{6} & NRc & Poprawiona prędkość obrotowa wytwornicy spalin \\
 & Nc & Prędkość obrotowa wytwornicy spalin \\
\cline{1-3}
\bottomrule
\end{tabular}

}
\end{table}

\begin{figure}[H]
    \centering
    \caption{\textbf{Istotność cech (Feature Importance) wyznaczona modelem XGBoost}\\ Oś pionowa przedstawia względną istotność cech}
    \includegraphics[width=0.8\linewidth]{Output/EDA/xgb_importance.png}
    \caption*{Źródło: opracowanie własne}
    \label{fig:xgb_importance}
\end{figure}

W celu optymalizacji czasu trenowania sieci zmniejszono wymiarowość danych. Na podstawie dendrogramu wybrano po jednej reprezentatywnej cesze dla każdej z grup. Pierwotnym kryterium była korelacja z~RUL. Otrzymano w~ten sposób zestaw:
\begin{verbatim}
    Operational setting 1, Operational setting 2, Ps30, P15, phi, Nc
\end{verbatim}

Ps30, Nc oraz phi znajdują się w~zbiorze sześciu najważniejszych cech ze względu na RUL wyłonionych w~analizie XGBoost.

Przeanalizowano rozkłady Operational setting 2 oraz P15 i~usunięto je ze zbioru, ponieważ uznano je za zmienne kategoryczne, a~nie ciągłe w~przeciwieństwie do pozostałych z~zestawu. Po analizie macierzy korelacji wybrano T30 oraz W31 na ich zastępstwo, otrzymując następujący zestaw cech:
\begin{verbatim}
    Operational setting 1, T30, Ps30, W31, phi, Nc
\end{verbatim}

Z macierzy korelacji Pearsona, w~której znajdują się tylko wybrane cechy, wynika, że na 15 możliwych par 6 z~nich ma korelację powyżej (co do wartości bezwzględnej) 0.6, co może prowadzić do współliniowości.

Powyższe cechy mają następujące znaczenie:
\begin{itemize}
    \item \textbf{Ps30} Istotny wskaźnik sprawności sprężania - jednego z~trzech etapów pracy silnika lotniczego. Jego spadek to sygnał zużycia łopatek lub nieszczelności.
    \item \textbf{T30} Przyrost temperatury może świadczyć o~erozji łopatek, których spadek sprawności jest korygowany przez komputer pokładowy zwiększając prędkość obrotową silnika. Wyższa temperatura zwiększa też obciążenie cieplne całej konstrukcji. W~skrócie jest to wskaźnik strat energii.
    \item \textbf{Nc} Prędkość obrotowa głównego wału. Do utrzymania ciągu przy postępującej degradacji, konieczne jest zwiększenie prędkości obrotowej.
    \item \textbf{phi} Opisuje zapotrzebowanie silnika na paliwo w~stosunku do wytwarzanego ciśnienia. Wskaźnik sprawności spalania.
    \item \textbf{W31} Przepływ powietrza chłodzącego turbinę. Zmiany wskazują na degradację uszczelnień i~upływy energii.
    \item \textbf{Op. Set 1} Parametr zewnętrzny (nieopisany), który definiuje warunki brzegowe pracy wszystkich powyższych sensorów.
\end{itemize}

Wykonano wykresy punktowe czterech najbardziej skorelowanych cech z~RUL w~funkcji przelatanych cykli (rysunek~\ref{fig:scatter}). Kolory punktów odpowiadają liczbie pozostałych cykli do przeżycia: im kolor ciemniejszy tym bliżej do remontu silnika. 
W przypadku \textbf{Ps30} (ciśnienie statyczne na wylocie turbiny wysokiego ciśnienia) ciśnienie na początku życia silników wynosi od około 47 do 47,75 psia i~rośnie około 2\% do 48 - 48,5 psia. Wzrost ciśnienia statycznego wiąże się ze spowolnieniem przepływu powietrza w~sprężarce co jest spójne z~narracją postępującej eksploatacji i~deterioracji komponentów silnika.
Również \textbf{T50} (temperatura całkowita na wylocie turbiny niskiego ciśnienia) rośnie wraz z cyklami. Wraz ze spadkiem prędkości powietrza w~sprężarkach, powietrze, które jest wykorzystywane do chłodzenia komory spalania, mniej skutecznie odbiera ciepło. W~innych modelach przewiduje się, że temperatura rośnie w~wyniku deterioracji turbiny, co przekłada się na gorsze odebranie energii, jednak w~przypadku tego zestawu danych modelowany jest jedynie spadek sprawności sprężarki wysokiego ciśnienia. Zmiany temperatur wahają się od $1390-1420$ do $1420-1440$ (około $2\%$)
Stosunek wydatku paliwa do ciśnienia statycznego na wylocie Sprężarki Wysokiego Ciśnienia (\textbf{phi}) w~funkcji przelatanych cykli jest funkcją malejącą, ponieważ wydatek paliwa jest stały, a~ciśnienie statyczne rośnie.
Ciśnienie całkowite za sprężarką wysokiego ciśnienia (\textbf{P30}) maleje, ponieważ straty rosną (maleje jej sprawność) w~funkcji cykli.

Na każdym z~czterech wykresów można wyróżnić obszary, w~których dla tej samej wartości analizowanej cechy część silników posiada zapas cykli, podczas gdy istnieją silniki, które już wymagały remontu. Dodatkowo to podkreśla rysunek~\ref{fig:rul_vs_cycles}, gdzie RUL w~funkcji cykli jest funkcją liniową malejącą o~tym samym współczynniku nachylenia, lecz (w zależności od silnika) różnym współczynniku wolnym.

\begin{figure}[H]
    \centering
    \caption{\textbf{Wykresy punktowe dla 4 najbardziej skorelowanych cech z~RUL}}
    \includegraphics[width=0.8\linewidth]{Output/Scatterplots/RUL_vs_P30_scatterplot.png}
    \caption*{Źródło: opracowanie własne}
    \label{fig:scatter}
\end{figure}

\begin{figure}[H]
    \centering
    \caption{\textbf{Wykres RUL w~funkcji przelatanych cykli}}
    \includegraphics[width=0.8\linewidth]{Output/Scatterplots/RUL_vs_cycles_scatterplot.png}
    \caption*{Źródło: opracowanie własne}
    \label{fig:rul_vs_cycles}
\end{figure}

